%% Sections 1--2: Introduction and Background
%% To be included in main paper via %% Sections 1--2: Introduction and Background
%% To be included in main paper via %% Sections 1--2: Introduction and Background
%% To be included in main paper via %% Sections 1--2: Introduction and Background
%% To be included in main paper via \input{sections1_2}

\section{Introduction}
\label{sec:introduction}

Group Relative Policy Optimization \citep[GRPO;][]{deepseekr1} has emerged as the dominant reinforcement learning algorithm for large language model reasoning, powering systems such as DeepSeek-R1 \citep{deepseekr1}, Qwen \citep{qwen2025}, and their derivatives.
GRPO inherits the clipped surrogate objective from Proximal Policy Optimization \citep[PPO;][]{schulman2017ppo}---a trust-region mechanism designed to prevent destructive policy updates by bounding the importance sampling ratio between the current and data-generating policies.
This mechanism is widely regarded as essential to training stability.

Yet a striking pattern has emerged in the recent literature: at least six independent works have redesigned or removed this clipping mechanism, each from a different empirical motivation, and each reporting comparable or improved results.
DAPO \citep{dapo2025} introduces asymmetric clip bounds (\emph{Clip-Higher}) to mitigate entropy collapse.
GSPO \citep{gspo2025} replaces token-level importance sampling with sequence-level ratios, arguing the token-level formulation is ``fundamentally broken.''
CFPO \citep{cfpo2025} substitutes clipping with a quadratic TV-divergence penalty to eliminate dead-zone gradients.
ASPO \citep{aspo2025} inverts the importance ratio for positive-advantage tokens to correct gradient misallocation.
RGRA \citep{rgra2025} removes clipping entirely, retaining plain REINFORCE with group-relative advantages, and matches or exceeds GRPO performance.
\citet{chen2025iclr} analyze clipping as an implicit entropy regularizer and recommend its removal.
All of these modifications work.
The natural question is: \emph{why did each become necessary?}

\paragraph{Our answer.}
We provide a unifying diagnosis.
By instrumenting the PPO loss to track \texttt{pg\_clipfrac}---the fraction of tokens at which the clipped and unclipped objectives differ---we find that this quantity is \emph{exactly zero} across all standard GRPO training configurations we tested.
The clip operator in the PPO surrogate never binds: for every token, the clipped objective equals the unclipped objective, and the trust-region mechanism contributes nothing to the optimization.

The root cause is an architectural disconnection between the importance ratio and data staleness.
In classical PPO, the ratio denominator is the behavior policy that generated the training data.
In GRPO as implemented in standard frameworks (TRL, SLIME, veRL), the denominator uses log-probabilities \emph{recomputed} from the current model at each training step.
Since both numerator and denominator track the same policy, the ratio remains near unity regardless of how far the model has drifted from the policy that generated the rollout data.
We identify two independent causes: (1)~this recomputation of log-probabilities resets the ratio, and (2)~conservative learning rates ($10^{-6}$--$10^{-5}$) limit per-step policy change.
Either cause alone is sufficient to keep clipping inactive; together, they render it entirely inert.

\paragraph{Activating clipping: a controlled experiment.}
We introduce a \texttt{--use-rollout-logprobs} flag that replaces the ratio denominator with the stale rollout log-probabilities, reconnecting the importance ratio to the data-generating distribution.
Combined with tight clipping ($\varepsilon = 0.05$), this is the only configuration in our experimental grid that produces entropy \emph{increase} during training: entropy rises fourfold (from 0.28 to 1.19) while reward initially improves.
However, extended training over 310 rollouts reveals gradual reward decline---the model explores more diverse strategies but does not converge to better ones.
This is itself informative: it demonstrates that clipping is \emph{necessary but not sufficient} under genuine off-policyness, and that complementary mechanisms are needed for sustained improvement.

\paragraph{Contributions.}
Our contributions are:
\begin{enumerate}
    \item The first empirical measurement of clipping inactivity in GRPO training ($\texttt{pg\_clipfrac} = 0$) and identification of its architectural root cause: the PPO importance ratio uses recomputed log-probabilities that are disconnected from data staleness.

    \item The introduction of Effective Sample Size (ESS) as a principled offlineness metric for GRPO, providing a continuous measure of the divergence between the current policy and the data-generating distribution.

    \item A unifying diagnosis that explains why six independent recent works \citep{dapo2025,gspo2025,cfpo2025,aspo2025,rgra2025,chen2025iclr} each found it necessary to redesign or remove GRPO's clipping mechanism.

    \item A controlled factorial experiment showing what happens when clipping is activated: entropy increases and clipping engages as intended, but sustained reward improvement requires mechanisms beyond clipping alone.
\end{enumerate}

\paragraph{Relation to prior observations.}
Several prior works have noted symptoms of clipping inactivity without identifying the cause.
\citet{chen2025iclr} measure $\texttt{pg\_clipfrac} < 0.2\%$ on Qwen2.5-Math-7B at $\text{lr} = 5 \times 10^{-7}$ but attribute the low rate to conservative hyperparameters.
\citet{prosperity2025} report $\texttt{pg\_clipfrac} = 0.05\%$ at staleness $s = 0$ in their multi-policy framework but do not investigate the root cause.
\citet{grpo_secretly_prm} note that the importance ratio equals unity when the number of PPO epochs $\mu = 1$, and \citet{lambert2025} observes that clipping is ``irrelevant'' in this setting.
A TRL user reported $r_t = 1$ in issue \#2769, which the community dismissed as expected behavior.\footnote{\url{https://github.com/huggingface/trl/issues/2769}}
Our contribution is the causal diagnosis---the recomputation of log-probabilities in the PPO loss---and a controlled experiment demonstrating what changes when clipping is made functional.


\section{Background}
\label{sec:background}

\subsection{GRPO and the PPO Clipped Objective}
\label{sec:bg-grpo}

Group Relative Policy Optimization \citep{deepseekr1} generates multiple candidate responses $\{y_1, \dots, y_G\}$ for each prompt $x$ and computes \emph{group-relative advantages}: for each response $y_i$, the advantage $\hat{A}_i$ is derived from the reward $R(x, y_i)$ relative to the group statistics (mean and, optionally, standard deviation) of $\{R(x, y_j)\}_{j=1}^G$.
This eliminates the need for a learned value function, which is a primary source of instability and computational cost in standard PPO for language models.

The policy is then updated using PPO's clipped surrogate objective, applied at the token level.
For a token $a_t$ generated in context $s_t$, the loss is
\begin{equation}
\label{eq:ppo-clip}
\mathcal{L}^{\text{CLIP}}_t(\theta) = -\min\!\Big( r_t(\theta)\,\hat{A}_t,\; \text{clip}\big(r_t(\theta),\, 1{-}\varepsilon,\, 1{+}\varepsilon\big)\,\hat{A}_t \Big),
\end{equation}
where $\varepsilon$ is the clip threshold (typically $0.2$) and $r_t(\theta)$ is the importance sampling ratio
\begin{equation}
\label{eq:is-ratio}
r_t(\theta) = \frac{\pi_\theta(a_t \mid s_t)}{\pi_{\text{old}}(a_t \mid s_t)}.
\end{equation}
When $\hat{A}_t > 0$, the clip operator caps the ratio at $1 + \varepsilon$, preventing excessively large upward updates; when $\hat{A}_t < 0$, it floors the ratio at $1 - \varepsilon$, preventing the policy from moving too far away from dispreferred actions in a single step.
The fraction of tokens at which the clipped and unclipped terms in Equation~\ref{eq:ppo-clip} differ---denoted \texttt{pg\_clipfrac}---measures how often this trust-region constraint is active.

DAPO \citep{dapo2025} modifies this with asymmetric bounds ($\varepsilon_{\text{low}} = 0.2$, $\varepsilon_{\text{high}} = 0.28$), allowing the policy to move further toward preferred actions than away from dispreferred ones.
The standard GRPO training loop additionally includes a per-token KL penalty against a reference policy $\pi_{\text{ref}}$, which we hold fixed in all experiments.


\subsection{The Three Log-Probability References}
\label{sec:bg-three-logprobs}

Understanding why clipping is inert requires distinguishing three log-probability quantities that coexist in GRPO training.
We define them precisely, as conflation of these quantities is the source of the architectural disconnection we diagnose.

\begin{enumerate}
    \item \textbf{Rollout log-probabilities} ($\log \pi_{\text{rollout}}$): computed at data generation time, when the model produces candidate responses for each prompt.
    These represent the actual \emph{behavior policy}---the distribution from which the training data was sampled.
    In code, these are stored as \texttt{batch["rollout\_log\_probs"]}.

    \item \textbf{Old log-probabilities} ($\log \pi_{\text{old}}$): recomputed from the current model at the start of each training step (or PPO epoch).
    In code, these are stored as \texttt{batch["log\_probs"]}.
    In standard GRPO implementations, this quantity serves as the denominator of the importance ratio in Equation~\ref{eq:is-ratio}.

    \item \textbf{Current log-probabilities} ($\log \pi_\theta$): the model's output during the forward pass of the training step.
    This serves as the numerator of the importance ratio.
\end{enumerate}

The critical distinction is the identity of $\pi_{\text{old}}$ in Equation~\ref{eq:is-ratio}.
In classical PPO \citep{schulman2017ppo}, $\pi_{\text{old}} = \pi_{\text{rollout}}$: the ratio denominator is the behavior policy that generated the data, and the ratio measures how far the current policy has moved from the data-generating distribution.
In GRPO as implemented in standard frameworks---including TRL,\footnote{\url{https://github.com/huggingface/trl}} veRL \citep{verl2024}, and SLIME---$\pi_{\text{old}} = \pi_{\theta_{\text{step-start}}}$: the log-probabilities are recomputed from the current model parameters at the beginning of each training step.
This means both the numerator ($\pi_\theta$) and denominator ($\pi_{\text{old}}$) of $r_t(\theta)$ track the same underlying policy, differing only by the parameter updates accumulated during a single gradient step.
The importance ratio is therefore \emph{architecturally near unity}, regardless of how far $\pi_\theta$ has drifted from the rollout policy $\pi_{\text{rollout}}$.

This substitution disconnects the PPO trust region from data staleness.
The clipped objective in Equation~\ref{eq:ppo-clip} constrains how far $\pi_\theta$ moves from $\pi_{\text{old}}$---but since $\pi_{\text{old}}$ is reset to $\pi_\theta$ at each step, the constraint is vacuous.
The policy may drift arbitrarily far from the data-generating distribution $\pi_{\text{rollout}}$ without the clipping mechanism ever detecting or constraining this drift.


\subsection{Offlineness Metrics}
\label{sec:bg-offlineness}

To quantify the divergence between the current policy and the data-generating distribution, we employ the \textbf{Effective Sample Size} (ESS), a standard diagnostic in importance sampling \citep{kong1992ess,liu2001mc}:
\begin{equation}
\label{eq:ess}
\text{ESS} = \frac{\big(\sum_{i} w_i\big)^2}{\sum_{i} w_i^2}, \quad w_i = \frac{\pi_\theta(a_i \mid s_i)}{\pi_{\text{rollout}}(a_i \mid s_i)},
\end{equation}
where the weights $w_i$ are the importance ratios computed against the \emph{rollout} policy (not the recomputed $\pi_{\text{old}}$).
Normalized to the interval $[0, 1]$, $\text{ESS} \approx 1.0$ indicates that the current policy closely matches the data-generating distribution (near on-policy training), while $\text{ESS} \ll 1.0$ indicates severe off-policyness---the training data is unrepresentative of the current policy, and importance-weighted estimates will have high variance.

We additionally track three complementary metrics:
\begin{itemize}
    \item \textbf{Mismatch KL} ($D_{\chi^2}$): a chi-squared divergence proxy measuring distributional mismatch between $\pi_\theta$ and $\pi_{\text{rollout}}$.
    \item \textbf{Clip fraction} (\texttt{clip\_frac\_02}): the fraction of per-token importance weights $w_i$ falling outside the interval $[0.8, 1.2]$, providing a threshold-based staleness indicator.
    \item \textbf{Maximum ratio} (\texttt{max\_ratio}): the largest importance weight $\max_i w_i$, capturing extreme off-policy tokens that may destabilize training.
\end{itemize}

These metrics are computed using rollout log-probabilities as the reference, ensuring they reflect actual data staleness rather than the artificial near-unity ratios produced by the recomputed $\pi_{\text{old}}$.
The distinction between ESS (which uses $\pi_{\text{rollout}}$) and \texttt{pg\_clipfrac} (which uses recomputed $\pi_{\text{old}}$) is central to our diagnosis: a training run can exhibit $\text{ESS} = 0.988$ (measurable off-policyness) while simultaneously showing $\texttt{pg\_clipfrac} = 0$ (no clipping activity), precisely because the two metrics reference different denominators.


\section{Introduction}
\label{sec:introduction}

Group Relative Policy Optimization \citep[GRPO;][]{deepseekr1} has emerged as the dominant reinforcement learning algorithm for large language model reasoning, powering systems such as DeepSeek-R1 \citep{deepseekr1}, Qwen \citep{qwen2025}, and their derivatives.
GRPO inherits the clipped surrogate objective from Proximal Policy Optimization \citep[PPO;][]{schulman2017ppo}---a trust-region mechanism designed to prevent destructive policy updates by bounding the importance sampling ratio between the current and data-generating policies.
This mechanism is widely regarded as essential to training stability.

Yet a striking pattern has emerged in the recent literature: at least six independent works have redesigned or removed this clipping mechanism, each from a different empirical motivation, and each reporting comparable or improved results.
DAPO \citep{dapo2025} introduces asymmetric clip bounds (\emph{Clip-Higher}) to mitigate entropy collapse.
GSPO \citep{gspo2025} replaces token-level importance sampling with sequence-level ratios, arguing the token-level formulation is ``fundamentally broken.''
CFPO \citep{cfpo2025} substitutes clipping with a quadratic TV-divergence penalty to eliminate dead-zone gradients.
ASPO \citep{aspo2025} inverts the importance ratio for positive-advantage tokens to correct gradient misallocation.
RGRA \citep{rgra2025} removes clipping entirely, retaining plain REINFORCE with group-relative advantages, and matches or exceeds GRPO performance.
\citet{chen2025iclr} analyze clipping as an implicit entropy regularizer and recommend its removal.
All of these modifications work.
The natural question is: \emph{why did each become necessary?}

\paragraph{Our answer.}
We provide a unifying diagnosis.
By instrumenting the PPO loss to track \texttt{pg\_clipfrac}---the fraction of tokens at which the clipped and unclipped objectives differ---we find that this quantity is \emph{exactly zero} across all standard GRPO training configurations we tested.
The clip operator in the PPO surrogate never binds: for every token, the clipped objective equals the unclipped objective, and the trust-region mechanism contributes nothing to the optimization.

The root cause is an architectural disconnection between the importance ratio and data staleness.
In classical PPO, the ratio denominator is the behavior policy that generated the training data.
In GRPO as implemented in standard frameworks (TRL, SLIME, veRL), the denominator uses log-probabilities \emph{recomputed} from the current model at each training step.
Since both numerator and denominator track the same policy, the ratio remains near unity regardless of how far the model has drifted from the policy that generated the rollout data.
We identify two independent causes: (1)~this recomputation of log-probabilities resets the ratio, and (2)~conservative learning rates ($10^{-6}$--$10^{-5}$) limit per-step policy change.
Either cause alone is sufficient to keep clipping inactive; together, they render it entirely inert.

\paragraph{Activating clipping: a controlled experiment.}
We introduce a \texttt{--use-rollout-logprobs} flag that replaces the ratio denominator with the stale rollout log-probabilities, reconnecting the importance ratio to the data-generating distribution.
Combined with tight clipping ($\varepsilon = 0.05$), this is the only configuration in our experimental grid that produces entropy \emph{increase} during training: entropy rises fourfold (from 0.28 to 1.19) while reward initially improves.
However, extended training over 310 rollouts reveals gradual reward decline---the model explores more diverse strategies but does not converge to better ones.
This is itself informative: it demonstrates that clipping is \emph{necessary but not sufficient} under genuine off-policyness, and that complementary mechanisms are needed for sustained improvement.

\paragraph{Contributions.}
Our contributions are:
\begin{enumerate}
    \item The first empirical measurement of clipping inactivity in GRPO training ($\texttt{pg\_clipfrac} = 0$) and identification of its architectural root cause: the PPO importance ratio uses recomputed log-probabilities that are disconnected from data staleness.

    \item The introduction of Effective Sample Size (ESS) as a principled offlineness metric for GRPO, providing a continuous measure of the divergence between the current policy and the data-generating distribution.

    \item A unifying diagnosis that explains why six independent recent works \citep{dapo2025,gspo2025,cfpo2025,aspo2025,rgra2025,chen2025iclr} each found it necessary to redesign or remove GRPO's clipping mechanism.

    \item A controlled factorial experiment showing what happens when clipping is activated: entropy increases and clipping engages as intended, but sustained reward improvement requires mechanisms beyond clipping alone.
\end{enumerate}

\paragraph{Relation to prior observations.}
Several prior works have noted symptoms of clipping inactivity without identifying the cause.
\citet{chen2025iclr} measure $\texttt{pg\_clipfrac} < 0.2\%$ on Qwen2.5-Math-7B at $\text{lr} = 5 \times 10^{-7}$ but attribute the low rate to conservative hyperparameters.
\citet{prosperity2025} report $\texttt{pg\_clipfrac} = 0.05\%$ at staleness $s = 0$ in their multi-policy framework but do not investigate the root cause.
\citet{grpo_secretly_prm} note that the importance ratio equals unity when the number of PPO epochs $\mu = 1$, and \citet{lambert2025} observes that clipping is ``irrelevant'' in this setting.
A TRL user reported $r_t = 1$ in issue \#2769, which the community dismissed as expected behavior.\footnote{\url{https://github.com/huggingface/trl/issues/2769}}
Our contribution is the causal diagnosis---the recomputation of log-probabilities in the PPO loss---and a controlled experiment demonstrating what changes when clipping is made functional.


\section{Background}
\label{sec:background}

\subsection{GRPO and the PPO Clipped Objective}
\label{sec:bg-grpo}

Group Relative Policy Optimization \citep{deepseekr1} generates multiple candidate responses $\{y_1, \dots, y_G\}$ for each prompt $x$ and computes \emph{group-relative advantages}: for each response $y_i$, the advantage $\hat{A}_i$ is derived from the reward $R(x, y_i)$ relative to the group statistics (mean and, optionally, standard deviation) of $\{R(x, y_j)\}_{j=1}^G$.
This eliminates the need for a learned value function, which is a primary source of instability and computational cost in standard PPO for language models.

The policy is then updated using PPO's clipped surrogate objective, applied at the token level.
For a token $a_t$ generated in context $s_t$, the loss is
\begin{equation}
\label{eq:ppo-clip}
\mathcal{L}^{\text{CLIP}}_t(\theta) = -\min\!\Big( r_t(\theta)\,\hat{A}_t,\; \text{clip}\big(r_t(\theta),\, 1{-}\varepsilon,\, 1{+}\varepsilon\big)\,\hat{A}_t \Big),
\end{equation}
where $\varepsilon$ is the clip threshold (typically $0.2$) and $r_t(\theta)$ is the importance sampling ratio
\begin{equation}
\label{eq:is-ratio}
r_t(\theta) = \frac{\pi_\theta(a_t \mid s_t)}{\pi_{\text{old}}(a_t \mid s_t)}.
\end{equation}
When $\hat{A}_t > 0$, the clip operator caps the ratio at $1 + \varepsilon$, preventing excessively large upward updates; when $\hat{A}_t < 0$, it floors the ratio at $1 - \varepsilon$, preventing the policy from moving too far away from dispreferred actions in a single step.
The fraction of tokens at which the clipped and unclipped terms in Equation~\ref{eq:ppo-clip} differ---denoted \texttt{pg\_clipfrac}---measures how often this trust-region constraint is active.

DAPO \citep{dapo2025} modifies this with asymmetric bounds ($\varepsilon_{\text{low}} = 0.2$, $\varepsilon_{\text{high}} = 0.28$), allowing the policy to move further toward preferred actions than away from dispreferred ones.
The standard GRPO training loop additionally includes a per-token KL penalty against a reference policy $\pi_{\text{ref}}$, which we hold fixed in all experiments.


\subsection{The Three Log-Probability References}
\label{sec:bg-three-logprobs}

Understanding why clipping is inert requires distinguishing three log-probability quantities that coexist in GRPO training.
We define them precisely, as conflation of these quantities is the source of the architectural disconnection we diagnose.

\begin{enumerate}
    \item \textbf{Rollout log-probabilities} ($\log \pi_{\text{rollout}}$): computed at data generation time, when the model produces candidate responses for each prompt.
    These represent the actual \emph{behavior policy}---the distribution from which the training data was sampled.
    In code, these are stored as \texttt{batch["rollout\_log\_probs"]}.

    \item \textbf{Old log-probabilities} ($\log \pi_{\text{old}}$): recomputed from the current model at the start of each training step (or PPO epoch).
    In code, these are stored as \texttt{batch["log\_probs"]}.
    In standard GRPO implementations, this quantity serves as the denominator of the importance ratio in Equation~\ref{eq:is-ratio}.

    \item \textbf{Current log-probabilities} ($\log \pi_\theta$): the model's output during the forward pass of the training step.
    This serves as the numerator of the importance ratio.
\end{enumerate}

The critical distinction is the identity of $\pi_{\text{old}}$ in Equation~\ref{eq:is-ratio}.
In classical PPO \citep{schulman2017ppo}, $\pi_{\text{old}} = \pi_{\text{rollout}}$: the ratio denominator is the behavior policy that generated the data, and the ratio measures how far the current policy has moved from the data-generating distribution.
In GRPO as implemented in standard frameworks---including TRL,\footnote{\url{https://github.com/huggingface/trl}} veRL \citep{verl2024}, and SLIME---$\pi_{\text{old}} = \pi_{\theta_{\text{step-start}}}$: the log-probabilities are recomputed from the current model parameters at the beginning of each training step.
This means both the numerator ($\pi_\theta$) and denominator ($\pi_{\text{old}}$) of $r_t(\theta)$ track the same underlying policy, differing only by the parameter updates accumulated during a single gradient step.
The importance ratio is therefore \emph{architecturally near unity}, regardless of how far $\pi_\theta$ has drifted from the rollout policy $\pi_{\text{rollout}}$.

This substitution disconnects the PPO trust region from data staleness.
The clipped objective in Equation~\ref{eq:ppo-clip} constrains how far $\pi_\theta$ moves from $\pi_{\text{old}}$---but since $\pi_{\text{old}}$ is reset to $\pi_\theta$ at each step, the constraint is vacuous.
The policy may drift arbitrarily far from the data-generating distribution $\pi_{\text{rollout}}$ without the clipping mechanism ever detecting or constraining this drift.


\subsection{Offlineness Metrics}
\label{sec:bg-offlineness}

To quantify the divergence between the current policy and the data-generating distribution, we employ the \textbf{Effective Sample Size} (ESS), a standard diagnostic in importance sampling \citep{kong1992ess,liu2001mc}:
\begin{equation}
\label{eq:ess}
\text{ESS} = \frac{\big(\sum_{i} w_i\big)^2}{\sum_{i} w_i^2}, \quad w_i = \frac{\pi_\theta(a_i \mid s_i)}{\pi_{\text{rollout}}(a_i \mid s_i)},
\end{equation}
where the weights $w_i$ are the importance ratios computed against the \emph{rollout} policy (not the recomputed $\pi_{\text{old}}$).
Normalized to the interval $[0, 1]$, $\text{ESS} \approx 1.0$ indicates that the current policy closely matches the data-generating distribution (near on-policy training), while $\text{ESS} \ll 1.0$ indicates severe off-policyness---the training data is unrepresentative of the current policy, and importance-weighted estimates will have high variance.

We additionally track three complementary metrics:
\begin{itemize}
    \item \textbf{Mismatch KL} ($D_{\chi^2}$): a chi-squared divergence proxy measuring distributional mismatch between $\pi_\theta$ and $\pi_{\text{rollout}}$.
    \item \textbf{Clip fraction} (\texttt{clip\_frac\_02}): the fraction of per-token importance weights $w_i$ falling outside the interval $[0.8, 1.2]$, providing a threshold-based staleness indicator.
    \item \textbf{Maximum ratio} (\texttt{max\_ratio}): the largest importance weight $\max_i w_i$, capturing extreme off-policy tokens that may destabilize training.
\end{itemize}

These metrics are computed using rollout log-probabilities as the reference, ensuring they reflect actual data staleness rather than the artificial near-unity ratios produced by the recomputed $\pi_{\text{old}}$.
The distinction between ESS (which uses $\pi_{\text{rollout}}$) and \texttt{pg\_clipfrac} (which uses recomputed $\pi_{\text{old}}$) is central to our diagnosis: a training run can exhibit $\text{ESS} = 0.988$ (measurable off-policyness) while simultaneously showing $\texttt{pg\_clipfrac} = 0$ (no clipping activity), precisely because the two metrics reference different denominators.


\section{Introduction}
\label{sec:introduction}

Group Relative Policy Optimization \citep[GRPO;][]{deepseekr1} has emerged as the dominant reinforcement learning algorithm for large language model reasoning, powering systems such as DeepSeek-R1 \citep{deepseekr1}, Qwen \citep{qwen2025}, and their derivatives.
GRPO inherits the clipped surrogate objective from Proximal Policy Optimization \citep[PPO;][]{schulman2017ppo}---a trust-region mechanism designed to prevent destructive policy updates by bounding the importance sampling ratio between the current and data-generating policies.
This mechanism is widely regarded as essential to training stability.

Yet a striking pattern has emerged in the recent literature: at least six independent works have redesigned or removed this clipping mechanism, each from a different empirical motivation, and each reporting comparable or improved results.
DAPO \citep{dapo2025} introduces asymmetric clip bounds (\emph{Clip-Higher}) to mitigate entropy collapse.
GSPO \citep{gspo2025} replaces token-level importance sampling with sequence-level ratios, arguing the token-level formulation is ``fundamentally broken.''
CFPO \citep{cfpo2025} substitutes clipping with a quadratic TV-divergence penalty to eliminate dead-zone gradients.
ASPO \citep{aspo2025} inverts the importance ratio for positive-advantage tokens to correct gradient misallocation.
RGRA \citep{rgra2025} removes clipping entirely, retaining plain REINFORCE with group-relative advantages, and matches or exceeds GRPO performance.
\citet{chen2025iclr} analyze clipping as an implicit entropy regularizer and recommend its removal.
All of these modifications work.
The natural question is: \emph{why did each become necessary?}

\paragraph{Our answer.}
We provide a unifying diagnosis.
By instrumenting the PPO loss to track \texttt{pg\_clipfrac}---the fraction of tokens at which the clipped and unclipped objectives differ---we find that this quantity is \emph{exactly zero} across all standard GRPO training configurations we tested.
The clip operator in the PPO surrogate never binds: for every token, the clipped objective equals the unclipped objective, and the trust-region mechanism contributes nothing to the optimization.

The root cause is an architectural disconnection between the importance ratio and data staleness.
In classical PPO, the ratio denominator is the behavior policy that generated the training data.
In GRPO as implemented in standard frameworks (TRL, SLIME, veRL), the denominator uses log-probabilities \emph{recomputed} from the current model at each training step.
Since both numerator and denominator track the same policy, the ratio remains near unity regardless of how far the model has drifted from the policy that generated the rollout data.
We identify two independent causes: (1)~this recomputation of log-probabilities resets the ratio, and (2)~conservative learning rates ($10^{-6}$--$10^{-5}$) limit per-step policy change.
Either cause alone is sufficient to keep clipping inactive; together, they render it entirely inert.

\paragraph{Activating clipping: a controlled experiment.}
We introduce a \texttt{--use-rollout-logprobs} flag that replaces the ratio denominator with the stale rollout log-probabilities, reconnecting the importance ratio to the data-generating distribution.
Combined with tight clipping ($\varepsilon = 0.05$), this is the only configuration in our experimental grid that produces entropy \emph{increase} during training: entropy rises fourfold (from 0.28 to 1.19) while reward initially improves.
However, extended training over 310 rollouts reveals gradual reward decline---the model explores more diverse strategies but does not converge to better ones.
This is itself informative: it demonstrates that clipping is \emph{necessary but not sufficient} under genuine off-policyness, and that complementary mechanisms are needed for sustained improvement.

\paragraph{Contributions.}
Our contributions are:
\begin{enumerate}
    \item The first empirical measurement of clipping inactivity in GRPO training ($\texttt{pg\_clipfrac} = 0$) and identification of its architectural root cause: the PPO importance ratio uses recomputed log-probabilities that are disconnected from data staleness.

    \item The introduction of Effective Sample Size (ESS) as a principled offlineness metric for GRPO, providing a continuous measure of the divergence between the current policy and the data-generating distribution.

    \item A unifying diagnosis that explains why six independent recent works \citep{dapo2025,gspo2025,cfpo2025,aspo2025,rgra2025,chen2025iclr} each found it necessary to redesign or remove GRPO's clipping mechanism.

    \item A controlled factorial experiment showing what happens when clipping is activated: entropy increases and clipping engages as intended, but sustained reward improvement requires mechanisms beyond clipping alone.
\end{enumerate}

\paragraph{Relation to prior observations.}
Several prior works have noted symptoms of clipping inactivity without identifying the cause.
\citet{chen2025iclr} measure $\texttt{pg\_clipfrac} < 0.2\%$ on Qwen2.5-Math-7B at $\text{lr} = 5 \times 10^{-7}$ but attribute the low rate to conservative hyperparameters.
\citet{prosperity2025} report $\texttt{pg\_clipfrac} = 0.05\%$ at staleness $s = 0$ in their multi-policy framework but do not investigate the root cause.
\citet{grpo_secretly_prm} note that the importance ratio equals unity when the number of PPO epochs $\mu = 1$, and \citet{lambert2025} observes that clipping is ``irrelevant'' in this setting.
A TRL user reported $r_t = 1$ in issue \#2769, which the community dismissed as expected behavior.\footnote{\url{https://github.com/huggingface/trl/issues/2769}}
Our contribution is the causal diagnosis---the recomputation of log-probabilities in the PPO loss---and a controlled experiment demonstrating what changes when clipping is made functional.


\section{Background}
\label{sec:background}

\subsection{GRPO and the PPO Clipped Objective}
\label{sec:bg-grpo}

Group Relative Policy Optimization \citep{deepseekr1} generates multiple candidate responses $\{y_1, \dots, y_G\}$ for each prompt $x$ and computes \emph{group-relative advantages}: for each response $y_i$, the advantage $\hat{A}_i$ is derived from the reward $R(x, y_i)$ relative to the group statistics (mean and, optionally, standard deviation) of $\{R(x, y_j)\}_{j=1}^G$.
This eliminates the need for a learned value function, which is a primary source of instability and computational cost in standard PPO for language models.

The policy is then updated using PPO's clipped surrogate objective, applied at the token level.
For a token $a_t$ generated in context $s_t$, the loss is
\begin{equation}
\label{eq:ppo-clip}
\mathcal{L}^{\text{CLIP}}_t(\theta) = -\min\!\Big( r_t(\theta)\,\hat{A}_t,\; \text{clip}\big(r_t(\theta),\, 1{-}\varepsilon,\, 1{+}\varepsilon\big)\,\hat{A}_t \Big),
\end{equation}
where $\varepsilon$ is the clip threshold (typically $0.2$) and $r_t(\theta)$ is the importance sampling ratio
\begin{equation}
\label{eq:is-ratio}
r_t(\theta) = \frac{\pi_\theta(a_t \mid s_t)}{\pi_{\text{old}}(a_t \mid s_t)}.
\end{equation}
When $\hat{A}_t > 0$, the clip operator caps the ratio at $1 + \varepsilon$, preventing excessively large upward updates; when $\hat{A}_t < 0$, it floors the ratio at $1 - \varepsilon$, preventing the policy from moving too far away from dispreferred actions in a single step.
The fraction of tokens at which the clipped and unclipped terms in Equation~\ref{eq:ppo-clip} differ---denoted \texttt{pg\_clipfrac}---measures how often this trust-region constraint is active.

DAPO \citep{dapo2025} modifies this with asymmetric bounds ($\varepsilon_{\text{low}} = 0.2$, $\varepsilon_{\text{high}} = 0.28$), allowing the policy to move further toward preferred actions than away from dispreferred ones.
The standard GRPO training loop additionally includes a per-token KL penalty against a reference policy $\pi_{\text{ref}}$, which we hold fixed in all experiments.


\subsection{The Three Log-Probability References}
\label{sec:bg-three-logprobs}

Understanding why clipping is inert requires distinguishing three log-probability quantities that coexist in GRPO training.
We define them precisely, as conflation of these quantities is the source of the architectural disconnection we diagnose.

\begin{enumerate}
    \item \textbf{Rollout log-probabilities} ($\log \pi_{\text{rollout}}$): computed at data generation time, when the model produces candidate responses for each prompt.
    These represent the actual \emph{behavior policy}---the distribution from which the training data was sampled.
    In code, these are stored as \texttt{batch["rollout\_log\_probs"]}.

    \item \textbf{Old log-probabilities} ($\log \pi_{\text{old}}$): recomputed from the current model at the start of each training step (or PPO epoch).
    In code, these are stored as \texttt{batch["log\_probs"]}.
    In standard GRPO implementations, this quantity serves as the denominator of the importance ratio in Equation~\ref{eq:is-ratio}.

    \item \textbf{Current log-probabilities} ($\log \pi_\theta$): the model's output during the forward pass of the training step.
    This serves as the numerator of the importance ratio.
\end{enumerate}

The critical distinction is the identity of $\pi_{\text{old}}$ in Equation~\ref{eq:is-ratio}.
In classical PPO \citep{schulman2017ppo}, $\pi_{\text{old}} = \pi_{\text{rollout}}$: the ratio denominator is the behavior policy that generated the data, and the ratio measures how far the current policy has moved from the data-generating distribution.
In GRPO as implemented in standard frameworks---including TRL,\footnote{\url{https://github.com/huggingface/trl}} veRL \citep{verl2024}, and SLIME---$\pi_{\text{old}} = \pi_{\theta_{\text{step-start}}}$: the log-probabilities are recomputed from the current model parameters at the beginning of each training step.
This means both the numerator ($\pi_\theta$) and denominator ($\pi_{\text{old}}$) of $r_t(\theta)$ track the same underlying policy, differing only by the parameter updates accumulated during a single gradient step.
The importance ratio is therefore \emph{architecturally near unity}, regardless of how far $\pi_\theta$ has drifted from the rollout policy $\pi_{\text{rollout}}$.

This substitution disconnects the PPO trust region from data staleness.
The clipped objective in Equation~\ref{eq:ppo-clip} constrains how far $\pi_\theta$ moves from $\pi_{\text{old}}$---but since $\pi_{\text{old}}$ is reset to $\pi_\theta$ at each step, the constraint is vacuous.
The policy may drift arbitrarily far from the data-generating distribution $\pi_{\text{rollout}}$ without the clipping mechanism ever detecting or constraining this drift.


\subsection{Offlineness Metrics}
\label{sec:bg-offlineness}

To quantify the divergence between the current policy and the data-generating distribution, we employ the \textbf{Effective Sample Size} (ESS), a standard diagnostic in importance sampling \citep{kong1992ess,liu2001mc}:
\begin{equation}
\label{eq:ess}
\text{ESS} = \frac{\big(\sum_{i} w_i\big)^2}{\sum_{i} w_i^2}, \quad w_i = \frac{\pi_\theta(a_i \mid s_i)}{\pi_{\text{rollout}}(a_i \mid s_i)},
\end{equation}
where the weights $w_i$ are the importance ratios computed against the \emph{rollout} policy (not the recomputed $\pi_{\text{old}}$).
Normalized to the interval $[0, 1]$, $\text{ESS} \approx 1.0$ indicates that the current policy closely matches the data-generating distribution (near on-policy training), while $\text{ESS} \ll 1.0$ indicates severe off-policyness---the training data is unrepresentative of the current policy, and importance-weighted estimates will have high variance.

We additionally track three complementary metrics:
\begin{itemize}
    \item \textbf{Mismatch KL} ($D_{\chi^2}$): a chi-squared divergence proxy measuring distributional mismatch between $\pi_\theta$ and $\pi_{\text{rollout}}$.
    \item \textbf{Clip fraction} (\texttt{clip\_frac\_02}): the fraction of per-token importance weights $w_i$ falling outside the interval $[0.8, 1.2]$, providing a threshold-based staleness indicator.
    \item \textbf{Maximum ratio} (\texttt{max\_ratio}): the largest importance weight $\max_i w_i$, capturing extreme off-policy tokens that may destabilize training.
\end{itemize}

These metrics are computed using rollout log-probabilities as the reference, ensuring they reflect actual data staleness rather than the artificial near-unity ratios produced by the recomputed $\pi_{\text{old}}$.
The distinction between ESS (which uses $\pi_{\text{rollout}}$) and \texttt{pg\_clipfrac} (which uses recomputed $\pi_{\text{old}}$) is central to our diagnosis: a training run can exhibit $\text{ESS} = 0.988$ (measurable off-policyness) while simultaneously showing $\texttt{pg\_clipfrac} = 0$ (no clipping activity), precisely because the two metrics reference different denominators.


\section{Introduction}
\label{sec:introduction}

Group Relative Policy Optimization \citep[GRPO;][]{deepseekr1} has emerged as the dominant reinforcement learning algorithm for large language model reasoning, powering systems such as DeepSeek-R1 \citep{deepseekr1}, Qwen \citep{qwen2025}, and their derivatives.
GRPO inherits the clipped surrogate objective from Proximal Policy Optimization \citep[PPO;][]{schulman2017ppo}---a trust-region mechanism designed to prevent destructive policy updates by bounding the importance sampling ratio between the current and data-generating policies.
This mechanism is widely regarded as essential to training stability.

Yet a striking pattern has emerged in the recent literature: at least six independent works have redesigned or removed this clipping mechanism, each from a different empirical motivation, and each reporting comparable or improved results.
DAPO \citep{dapo2025} introduces asymmetric clip bounds (\emph{Clip-Higher}) to mitigate entropy collapse.
GSPO \citep{gspo2025} replaces token-level importance sampling with sequence-level ratios, arguing the token-level formulation is ``fundamentally broken.''
CFPO \citep{cfpo2025} substitutes clipping with a quadratic TV-divergence penalty to eliminate dead-zone gradients.
ASPO \citep{aspo2025} inverts the importance ratio for positive-advantage tokens to correct gradient misallocation.
RGRA \citep{rgra2025} removes clipping entirely, retaining plain REINFORCE with group-relative advantages, and matches or exceeds GRPO performance.
\citet{chen2025iclr} analyze clipping as an implicit entropy regularizer and recommend its removal.
All of these modifications work.
The natural question is: \emph{why did each become necessary?}

\paragraph{Our answer.}
We provide a unifying diagnosis.
By instrumenting the PPO loss to track \texttt{pg\_clipfrac}---the fraction of tokens at which the clipped and unclipped objectives differ---we find that this quantity is \emph{exactly zero} across all standard GRPO training configurations we tested.
The clip operator in the PPO surrogate never binds: for every token, the clipped objective equals the unclipped objective, and the trust-region mechanism contributes nothing to the optimization.

The root cause is an architectural disconnection between the importance ratio and data staleness.
In classical PPO, the ratio denominator is the behavior policy that generated the training data.
In GRPO as implemented in standard frameworks (TRL, SLIME, veRL), the denominator uses log-probabilities \emph{recomputed} from the current model at each training step.
Since both numerator and denominator track the same policy, the ratio remains near unity regardless of how far the model has drifted from the policy that generated the rollout data.
We identify two independent causes: (1)~this recomputation of log-probabilities resets the ratio, and (2)~conservative learning rates ($10^{-6}$--$10^{-5}$) limit per-step policy change.
Either cause alone is sufficient to keep clipping inactive; together, they render it entirely inert.

\paragraph{Activating clipping: a controlled experiment.}
We introduce a \texttt{--use-rollout-logprobs} flag that replaces the ratio denominator with the stale rollout log-probabilities, reconnecting the importance ratio to the data-generating distribution.
Combined with tight clipping ($\varepsilon = 0.05$), this is the only configuration in our experimental grid that produces entropy \emph{increase} during training: entropy rises fourfold (from 0.28 to 1.19) while reward initially improves.
However, extended training over 310 rollouts reveals gradual reward decline---the model explores more diverse strategies but does not converge to better ones.
This is itself informative: it demonstrates that clipping is \emph{necessary but not sufficient} under genuine off-policyness, and that complementary mechanisms are needed for sustained improvement.

\paragraph{Contributions.}
Our contributions are:
\begin{enumerate}
    \item The first empirical measurement of clipping inactivity in GRPO training ($\texttt{pg\_clipfrac} = 0$) and identification of its architectural root cause: the PPO importance ratio uses recomputed log-probabilities that are disconnected from data staleness.

    \item The introduction of Effective Sample Size (ESS) as a principled offlineness metric for GRPO, providing a continuous measure of the divergence between the current policy and the data-generating distribution.

    \item A unifying diagnosis that explains why six independent recent works \citep{dapo2025,gspo2025,cfpo2025,aspo2025,rgra2025,chen2025iclr} each found it necessary to redesign or remove GRPO's clipping mechanism.

    \item A controlled factorial experiment showing what happens when clipping is activated: entropy increases and clipping engages as intended, but sustained reward improvement requires mechanisms beyond clipping alone.
\end{enumerate}

\paragraph{Relation to prior observations.}
Several prior works have noted symptoms of clipping inactivity without identifying the cause.
\citet{chen2025iclr} measure $\texttt{pg\_clipfrac} < 0.2\%$ on Qwen2.5-Math-7B at $\text{lr} = 5 \times 10^{-7}$ but attribute the low rate to conservative hyperparameters.
\citet{prosperity2025} report $\texttt{pg\_clipfrac} = 0.05\%$ at staleness $s = 0$ in their multi-policy framework but do not investigate the root cause.
\citet{grpo_secretly_prm} note that the importance ratio equals unity when the number of PPO epochs $\mu = 1$, and \citet{lambert2025} observes that clipping is ``irrelevant'' in this setting.
A TRL user reported $r_t = 1$ in issue \#2769, which the community dismissed as expected behavior.\footnote{\url{https://github.com/huggingface/trl/issues/2769}}
Our contribution is the causal diagnosis---the recomputation of log-probabilities in the PPO loss---and a controlled experiment demonstrating what changes when clipping is made functional.


\section{Background}
\label{sec:background}

\subsection{GRPO and the PPO Clipped Objective}
\label{sec:bg-grpo}

Group Relative Policy Optimization \citep{deepseekr1} generates multiple candidate responses $\{y_1, \dots, y_G\}$ for each prompt $x$ and computes \emph{group-relative advantages}: for each response $y_i$, the advantage $\hat{A}_i$ is derived from the reward $R(x, y_i)$ relative to the group statistics (mean and, optionally, standard deviation) of $\{R(x, y_j)\}_{j=1}^G$.
This eliminates the need for a learned value function, which is a primary source of instability and computational cost in standard PPO for language models.

The policy is then updated using PPO's clipped surrogate objective, applied at the token level.
For a token $a_t$ generated in context $s_t$, the loss is
\begin{equation}
\label{eq:ppo-clip}
\mathcal{L}^{\text{CLIP}}_t(\theta) = -\min\!\Big( r_t(\theta)\,\hat{A}_t,\; \text{clip}\big(r_t(\theta),\, 1{-}\varepsilon,\, 1{+}\varepsilon\big)\,\hat{A}_t \Big),
\end{equation}
where $\varepsilon$ is the clip threshold (typically $0.2$) and $r_t(\theta)$ is the importance sampling ratio
\begin{equation}
\label{eq:is-ratio}
r_t(\theta) = \frac{\pi_\theta(a_t \mid s_t)}{\pi_{\text{old}}(a_t \mid s_t)}.
\end{equation}
When $\hat{A}_t > 0$, the clip operator caps the ratio at $1 + \varepsilon$, preventing excessively large upward updates; when $\hat{A}_t < 0$, it floors the ratio at $1 - \varepsilon$, preventing the policy from moving too far away from dispreferred actions in a single step.
The fraction of tokens at which the clipped and unclipped terms in Equation~\ref{eq:ppo-clip} differ---denoted \texttt{pg\_clipfrac}---measures how often this trust-region constraint is active.

DAPO \citep{dapo2025} modifies this with asymmetric bounds ($\varepsilon_{\text{low}} = 0.2$, $\varepsilon_{\text{high}} = 0.28$), allowing the policy to move further toward preferred actions than away from dispreferred ones.
The standard GRPO training loop additionally includes a per-token KL penalty against a reference policy $\pi_{\text{ref}}$, which we hold fixed in all experiments.


\subsection{The Three Log-Probability References}
\label{sec:bg-three-logprobs}

Understanding why clipping is inert requires distinguishing three log-probability quantities that coexist in GRPO training.
We define them precisely, as conflation of these quantities is the source of the architectural disconnection we diagnose.

\begin{enumerate}
    \item \textbf{Rollout log-probabilities} ($\log \pi_{\text{rollout}}$): computed at data generation time, when the model produces candidate responses for each prompt.
    These represent the actual \emph{behavior policy}---the distribution from which the training data was sampled.
    In code, these are stored as \texttt{batch["rollout\_log\_probs"]}.

    \item \textbf{Old log-probabilities} ($\log \pi_{\text{old}}$): recomputed from the current model at the start of each training step (or PPO epoch).
    In code, these are stored as \texttt{batch["log\_probs"]}.
    In standard GRPO implementations, this quantity serves as the denominator of the importance ratio in Equation~\ref{eq:is-ratio}.

    \item \textbf{Current log-probabilities} ($\log \pi_\theta$): the model's output during the forward pass of the training step.
    This serves as the numerator of the importance ratio.
\end{enumerate}

The critical distinction is the identity of $\pi_{\text{old}}$ in Equation~\ref{eq:is-ratio}.
In classical PPO \citep{schulman2017ppo}, $\pi_{\text{old}} = \pi_{\text{rollout}}$: the ratio denominator is the behavior policy that generated the data, and the ratio measures how far the current policy has moved from the data-generating distribution.
In GRPO as implemented in standard frameworks---including TRL,\footnote{\url{https://github.com/huggingface/trl}} veRL \citep{verl2024}, and SLIME---$\pi_{\text{old}} = \pi_{\theta_{\text{step-start}}}$: the log-probabilities are recomputed from the current model parameters at the beginning of each training step.
This means both the numerator ($\pi_\theta$) and denominator ($\pi_{\text{old}}$) of $r_t(\theta)$ track the same underlying policy, differing only by the parameter updates accumulated during a single gradient step.
The importance ratio is therefore \emph{architecturally near unity}, regardless of how far $\pi_\theta$ has drifted from the rollout policy $\pi_{\text{rollout}}$.

This substitution disconnects the PPO trust region from data staleness.
The clipped objective in Equation~\ref{eq:ppo-clip} constrains how far $\pi_\theta$ moves from $\pi_{\text{old}}$---but since $\pi_{\text{old}}$ is reset to $\pi_\theta$ at each step, the constraint is vacuous.
The policy may drift arbitrarily far from the data-generating distribution $\pi_{\text{rollout}}$ without the clipping mechanism ever detecting or constraining this drift.


\subsection{Offlineness Metrics}
\label{sec:bg-offlineness}

To quantify the divergence between the current policy and the data-generating distribution, we employ the \textbf{Effective Sample Size} (ESS), a standard diagnostic in importance sampling \citep{kong1992ess,liu2001mc}:
\begin{equation}
\label{eq:ess}
\text{ESS} = \frac{\big(\sum_{i} w_i\big)^2}{\sum_{i} w_i^2}, \quad w_i = \frac{\pi_\theta(a_i \mid s_i)}{\pi_{\text{rollout}}(a_i \mid s_i)},
\end{equation}
where the weights $w_i$ are the importance ratios computed against the \emph{rollout} policy (not the recomputed $\pi_{\text{old}}$).
Normalized to the interval $[0, 1]$, $\text{ESS} \approx 1.0$ indicates that the current policy closely matches the data-generating distribution (near on-policy training), while $\text{ESS} \ll 1.0$ indicates severe off-policyness---the training data is unrepresentative of the current policy, and importance-weighted estimates will have high variance.

We additionally track three complementary metrics:
\begin{itemize}
    \item \textbf{Mismatch KL} ($D_{\chi^2}$): a chi-squared divergence proxy measuring distributional mismatch between $\pi_\theta$ and $\pi_{\text{rollout}}$.
    \item \textbf{Clip fraction} (\texttt{clip\_frac\_02}): the fraction of per-token importance weights $w_i$ falling outside the interval $[0.8, 1.2]$, providing a threshold-based staleness indicator.
    \item \textbf{Maximum ratio} (\texttt{max\_ratio}): the largest importance weight $\max_i w_i$, capturing extreme off-policy tokens that may destabilize training.
\end{itemize}

These metrics are computed using rollout log-probabilities as the reference, ensuring they reflect actual data staleness rather than the artificial near-unity ratios produced by the recomputed $\pi_{\text{old}}$.
The distinction between ESS (which uses $\pi_{\text{rollout}}$) and \texttt{pg\_clipfrac} (which uses recomputed $\pi_{\text{old}}$) is central to our diagnosis: a training run can exhibit $\text{ESS} = 0.988$ (measurable off-policyness) while simultaneously showing $\texttt{pg\_clipfrac} = 0$ (no clipping activity), precisely because the two metrics reference different denominators.
